
\section{Introduction}

[2C-NEW-OPEN] This article determines the dimensions and almost-sure
frequencies of Hermite continued-fraction window constraints, the Hausdorff
dimensions of modular digit-residue sets, and the finite-state tail gauge
theory, along the selected convergent denominators of $e$.

Corollary~\ref{cor:liminf-gauges} classifies every nondecreasing gauge by
its lower limit along the recurrence scales, even when the normalized
cost oscillates. Theorem~\ref{thm:finite-state-tail} shows how this
classification changes under constraints on consecutive defects.
The largest Perron root is taken only over components in which a path
can revisit an extreme defect infinitely often.
[2C-NEW-NOT] We use the notation of the core paper (2A):
$S_k=\{x\in\R/\Z:\Lambda(x)=k/4\}$ for the exact levels and $d_j(x)$
for the consecutive-defect sequence; proofs of the quoted gauge results
are in 2A.
In Section~\ref{sec:windows-residues}, Theorems~\ref{thm:windows}--\ref{thm:no-short-dim}
determine the Hausdorff dimensions and almost-sure frequency of Hermite
continued-fraction window constraints, Theorem~\ref{thm:residues} proves
that eventual periodicity of ordinary continued-fraction digits modulo $4$ has
the same Hausdorff dimension as the fixed alphabet $1\bmod 4$, bounded strictly
between $357/512$ and $179/256$, Theorem~\ref{thm:residues-general} extends this
to every modulus $m\ge2$ (via a uniform $24$-Lipschitz sandwich, an unconditional
transfer-operator Hausdorff-dimension framework, and exact integer certificates
for $m\in\{2,3,4,5,6\}$), Theorem~\ref{thm:residues-asymptotic} proves the
global strict separation $\dimH F_{m,m}<\dimH F_{m,1}=\dimH E_m$ for all $m\ge2$
together with the exact two-term large-$m$ asymptotic expansion governed by the
reciprocal Fibonacci constant $\psi_{\mathrm{Fib}}=\sum_{k=1}^\infty 1/F_k$ and
its metallic-continuant generalizations $C_r=\sum_{k=1}^\infty 1/q_k(r)$, and
Theorem~\ref{thm:multi-residue-asymptotic} generalizes this to finite
multi-residue alphabets $F_{m,S,B}$---giving a self-contained Bowen-free proof
of the classical Poincar\'e-abscissa dimension formula for $E_S$, the universal phase transition
$\dimH F_{m,S,B}\to\max(1/2,s_0(S))$, the three-term Euler--Mascheroni
subcritical expansion to $O((\log m)^2/m^3)$, and the supercritical
$\Theta(m^{-2s_0(S)})$ convergence rate.

\subsection{Setup}

The continued fraction of $\ee$ is $[2;1,2,1,1,4,1,1,6,\dots]$. Let $P_j/Q_j$
be the convergent immediately preceding the partial quotient $2j$; the first
few are $3/1,19/7,193/71,2721/1001$. The denominators satisfy
\begin{equation}\label{eq:rec}
Q_{j+1}=(4j+2)Q_j+Q_{j-1}\quad(j\ge2),\qquad Q_1=1,\ Q_2=7,
\end{equation}
so $Q_{j+1}/Q_j\sim 4j$ and $\log Q_j\sim j\log j$ (super-exponential growth).
Write $\|y\|=\dist(y,\Z)$ and define
\[
\Lam(x):=\limsup_{j\to\infty} j\,\|Q_j x\|\in[0,\infty].
\]
More generally, for $\alpha\ge0$ put
$\Lam_\alpha(x):=\limsup_{j\to\infty} j^\alpha\|Q_j x\|\in[0,\infty]$,
so that $\Lam=\Lam_1$.

\subsection{Relation to prior work}\label{sec:prior}

For the modular continued-fraction alphabets $F_{m,r}=\{[0;a_1,a_2,\dots]:a_n\equiv r\pmod m\}$,
the underlying infinite-alphabet dimension theory is that of Mauldin and
Urba\'nski \cite{MU96} (see \cite[Theorem~2.3]{CLU}), Hensley \cite{Hensley1996}
(whose polynomial-time paper carries a published erratum), Jenkinson and
Pollicott \cite{JP2018}, and Chousionis, Leykekhman and Urba\'nski \cite{CLU},
who proved that every arithmetic-progression alphabet has full dimension
spectrum, while Chitanga, Lemmens and Nussbaum \cite{CLN2026} recently settled
full dimension spectra $[0,\dimH F_{m,r}]$ for single arithmetic progressions
(not addressing our uniform Lipschitz invariance across $E_m$, global strict
separation $\dimH F_{m,m}<\dimH F_{m,1}$, constants $C_r$, or multi-residue
asymptotics). Asymptotic expansions of the Hausdorff dimension when generators
with small derivative are added to or removed from a conformal IFS were
developed by Das, Fishman, Simmons and Urba\'nski \cite{DFSU}, who
strengthened Hensley's expansion \cite{Hensley92} for partial quotients $\le N$
and Good's estimates \cite{Good41} for partial quotients $\ge N$. The families
$F_{m,S,B}$ fall within the scope of their perturbation theorem
\cite[Theorem~4.1]{DFSU} (fixed finite $S$, added generators $\ge m-\max B$);
in their notation the leading constant is $1/\bar c=\sum_{n\ge0}\sum_{w\in S^n}q(w)^{-1}=C_S$
when $s_0(S)<1/2$. Arithmetic-progression alphabets are not treated in
\cite{DFSU}, and Theorems~\ref{thm:residues-asymptotic}
and~\ref{thm:multi-residue-asymptotic} make the first terms explicit for them;
the $\log m/m^2$ coefficient has the same structure as Hensley's
$72/\pi^4=2(6/\pi^2)^2$ (cf.\ \cite[(1.4)]{DFSU}). Part~(i) of
Theorem~\ref{thm:multi-residue-asymptotic} is the classical Bowen formula
written as the abscissa of a Poincar\'e series; our contribution there is a
proof that avoids the pressure formalism, together with its formalization.
The upper bound in part~(v), and hence the supercritical case of part~(ii),
follow from the Heinemann--Urba\'nski estimate \cite{HU02} (in the form of
\cite[Theorem~2.4]{CLU}) combined with the Lyapunov bound
\cite[Proposition~2.6]{CLU}, and the case $S=\emptyset$ of part~(ii) follows
from Good's theorem by monotonicity. Part~(iii) is a quantitative form of the
Laurent expansion of $\zeta$ at $1$ \cite[(25.2.4)]{DLMF}. A targeted negative
search (arXiv, Crossref and web search; coverage recorded in the verification
archive) found no prior appearance of $C_r$, of $\psi_{\mathrm{Fib}}$, or of the
coefficient $\frac12 d^2C_S(\gamma C_S-C_S')$ as asymptotic dimension constants
of residue-class digit sets; this does not exclude such an appearance.
% Theorem/equation numbers cited for CLU and DFSU above are those of arXiv:1805.11904v1 and
% arXiv:2007.10554v2; the numbering of the published versions has not been verified
% (verification/paper2_novelty_review_20261125.md).
In
Section~\ref{sec:multi-residue}, Theorem~\ref{thm:multi-residue-asymptotic}
extends this analysis to finite multi-residue alphabets $F_{m,S,B}$, giving a
self-contained Bowen-free proof of the classical Poincar\'e-abscissa formula
$\dimH E_S=s_0(S)$, the phase transition
$\lim_{m\to\infty}\dimH F_{m,S,B}=\max(1/2,s_0(S))$, the three-term
Euler--Mascheroni subcritical asymptotic expansion to $O((\log m)^2/m^3)$, and
the supercritical convergence rate $\Theta(m^{-2s_0(S)})$.
[2C-NEW-S30] The main new contributions are the window dimension and
short-window laws (Theorems~\ref{thm:windows}--\ref{thm:no-short-dim}),
the residue dimensions and strict separation (Theorems~\ref{thm:residues},
\ref{thm:residues-general}, and~\ref{thm:residues-asymptotic}(ii), with
Table~\ref{tab:residues-general}), explicit residue-class instances of known
asymptotic methods (Theorems~\ref{thm:residues-asymptotic}(iii)--(iv)
and~\ref{thm:multi-residue-asymptotic}, with Tables~\ref{tab:residues-asymptotic}
and~\ref{tab:multi-residue-asymptotic}), and
the finite-state tail theorem (Theorem~\ref{thm:finite-state-tail}).

